\documentclass[a4paper,12pt]{article}
\input{LaTeX-ToolKit/template.tex}
\title{}
\author{沈威宇}
\date{\temtoday}
\pagestyle{empty}
\begin{document}
\temdoc
\sct{First-order ODE}
\ssc{First-order ODE}
\[M(x,y)+N(x,y)\dv{y}{x}=0\]
\ssc{Separable equation}
\[\dv{y}{x}=g(x)h(y)\]
Family of solutions:
\[\int\frac{1}{h(y)}\dd{y}=\int g(x)\dd{x}+C\]
Singular solution:
\[h(y)=0\]
\ssc{First-order linear equation}
\[\dv{y}{x}+P(x)y=Q(x)\]
Integrating factor:
\[\mu(x)=e^{\int P(x)\dd{x}}\]
General solution:
\[y(x)=\frac{1}{\mu(x)}\qty(\int\mu(x)Q(x)\dd{x}+C)\]
\ssc{Exact equation}
On a simply-connected open set,
\[M(x,y)\dd{x}+N(x,y)\dd{y}=0\]
where $\frac{\partial^2M}{\partial x\partial y}$ and $\frac{\partial^2N}{\partial x\partial y}$ is continuous and
\[\pdv{M}{y}=\pdv{N}{x}\]
There exists a function $F$ such that
\[\pdv{F}{x}=M,\quad\pdv{F}{y}=N\]
and that $\frac{\partial^2F}{\partial x\partial y}$ is continuous, called the potential function.
\[\ba
F(x,y)&=\int M(x,y)\dd{x}+\int N(x,y)-\pdv{}{y}\qty(\int M(x,y)\dd{x})\dd{y}\\
      &=\int N(x,y)\dd{y}+\int M(x,y)-\pdv{}{x}\qty(\int N(x,y)\dd{y})\dd{x}
\ea\]
Implicit solution:
\[F(x,y)=C\]
\ssc{Integrating factors to make inexact equations exact}
On a simply-connected open set,
\[M(x,y)\dd{x}+N(x,y)\dd{y}=0\]
Integrating factor $\mu(x,y)$ such that
\[\mu(x,y)M(x,y)\dd{x}+\mu(x,y)N(x,y)\dd{y}=0\]
is exact.

If
\[\frac{\pdv{M}{y}-\pdv{N}{x}}{N}\]
depends only on $x$,
\[\mu(x)=e^{\int\frac{\pdv{M}{y}-\pdv{N}{x}}{N}\dd{x}}\]
If
\[\frac{\pdv{N}{x}-\pdv{M}{y}}{M}\]
depends only on $y$,
\[\mu(y)=e^{\int\frac{\pdv{N}{x}-\pdv{M}{y}}{M}\dd{x}}\]
\ssc{Homogeneous equation}
\[\dv{y}{x}=F(x,y)\]
\[\forall t\ne0\colon F(tx,ty)=F(x,y)\]
There exists a function $f$ such that
\[\forall x\ne0,(x,y)\in D_F\colon F(x,y)=f\qty(\frac{y}{x})\]
\[v=\frac{y}{x}\]
Solution:
\[\int\frac{\dd{v}}{f(v)-v}=\ln|x|+C\]
If $F$ is defined on $x\ne0$, also implicit solution
\[F(0,y)=0\]
\ssc{Bernoulli equation}
\[\dv{y}{x}+P(x)y=Q(x)y^n,\quad n\in\mathbb{R}\setminus\{0,1\}\]
\[v=y^{1-n}\]
\[\mu(x)=e^{\int (1-n)P(x)\dd{x}}\]
\[v(x)=\frac{1}{\mu(x)}\qty(\int\mu(x)(1-n)Q(x)\dd{x}+C)\]
\ssc{Riccati equation}
\[\dv{y}{x}=q_0(x)+q_1(x)y+q_2(x)y^2,\]
where $q_2(x)$ is nonzero and differentiable.

If $q_2(x)=1$, it is called canonical form.
\[v=yq_2\]
\[a(x)=q_0q_2\]
\[b(x)=q_1+\frac{q_2'}{q_2}\]
\[\dv{v}{x}=a(x)+b(x)v+v^2\]
\[u=\exp(-\int v\dd{x})\]
\[u''-bu'+au=0\]
\[y=-q_2^{-1}\qty(\ln|u|)'\]
\ssc{Slope being function of linear function}
\[\dv{y}{x}=f(u),\quad u=ax+by+c\]
\[\int\frac{1}{a+bf(u)}\dd{u}=x+C\]
\ssc{Slope being function of linear function over linear function}
\[\dv{y}{x}=F\qty(\frac{ax+by+c}{px+qy+r}),\quad aq-bp\neq 0.\]
We can find a unique solution of $ax+by+c=px+qy+r=0$, $x=\alpha$ and $y=\beta$. Let
\[f(t)=F\qty(\frac{a+bt}{p+qt})\]
\[\dv{y}{x}=f(u),\quad u=\frac{y-\beta}{x-\alpha}\]
Solutions:
\[\int\frac{1}{f(u)-u}\dd{u}=\ln|x-\alpha|+C\]
or
\[f(u)=u\]
\sct{Linear ODE}
\ssc{Basics}
\sssc{Linear ODE}
\[\sum_{i=0}^na_i(x)y^{(i)}=f(x),\]
in which $a_i(x)\colon\Omega_i\subseteq\bbR\to\bbR$ are called coefficients and $f(x)\colon\Omega\subseteq\bbR\to\bbR^m$ is called forcing term.
\sssc{Linear differential operator}
A linear differential operator is a polynomial of the differential operators $D$, i.e., $\dv{}{x}$ for independent variable $x$, and its order is the degree of the polynomial.
\sssc{Linearity}
Let $L$ be the linear differential operator $L[\mb{y}]=\sum_{i=0}^na_i(x)\mb{y}^{(i)}$. Then $L[\alpha\mb{y}_1+\beta\mb{y}_2]=\alpha L[\mb{y}_1]+\beta L[\mb{y}_2]$.
\sssc{Homogeneity}
A linear ODE is called homogeneous iff the forcing term is 0; otherwise, it is called non-homogeneous.
\sssc{Superposition principle and solutions of homogeneous linear ODE as vector space}
If $y_1(x)$ and $y_2(x)$ are both solutions to a homogeneous linear ODE on an interval, then $c_1y_1(x)+c_2y_2(x)$ is also a solution to it on the interval for any $c_1,c_2\in\bbR$. Thus, the general solutions to a homogeneous linear ODE is a vector space.
\sssc{Complementary or associated (homogeneous) equation, natural, complementary, or homogeneous solution or function, particularly solution, and solutions of non-homogeneous linear ODE as affine space}
Given an $n$th-order linear ODE
\[\sum_{i=0}^na_i(x)y^{(i)}=f(x),\]
\[\sum_{i=0}^na_i(x)y^{(i)}=0,\]
is called its complementary or associated (homogeneous) equation, and the general solutions of the original equation is the sum of the general solutions $y_c$ of the complementary equation, called natural, complementary, or homogeneous solution or function, and a particular solution $y_p$. Thus, the general solutions to a non-homogeneous linear ODE is an affine space whose associated vector space is the vector space of the solutions to its complementary equation.
\sssc{Wrońskian or Wronskian}
For $n$ functions $f_i\colon\Omega_i\subseteq\bbR\to\bbR$ for all $i\in\bbN_{\le n}$ that are $(n-1)$ times differentiable on an interval $I$, the Wrońskian (also called Wronskian) of them, $W(f_1,f_2,\ldots,f_n)$, is a function on $I$ defined by
\[W(f_1,f_2,\ldots,f_n)(x)=\begin{vmatrix}
f_1(x) & f_2(x) & \cdots & f_n(x) \\
f_1'(x) & f_2(x) & \cdots & f_n'(x) \\
\vdots & \vdots & \ddots & \vdots \\
f_1^{(n-1)}(x) & f_2^{(n-1)}(x) & \cdots & f_n^{(n-1)}(x)
\end{vmatrix}.\]
\sssc{Non-zero Wrońskian implies linear independence}
For $n$ functions $f_i\colon\Omega_i\subseteq\bbR\to\bbR$ for all $i\in\bbN_{\le n}$ that are $(n-1)$ times differentiable on an interval $I$, if $W(f_1,f_2,\ldots,f_n)\ne0$, then they are linearly independent.
\sssc{General Abel's (differential equation) identity/formula}
For $n$ $(n-1)$ times differentiable solutions $y_1,y_2,\ldots,y_n$ of $y^{(n)}+\sum_{i=0}^{n-1}p_i(x)y^{(i)}=0$ on an interval $I$ where $p_i(x)$ are continuous for all $i\in[0,n-1]\cap\bbZ$, for any $x_0\in I$,
\[W(y_1,y_2,\ldots,y_n)(x)=W(y_1,y_2,\ldots,y_n)(x_0)\exp(-\int_{x_0}^Xp_{n-1}(t)\dd{t}).\]
\sssc{Non-zero Wrońskian of solutions of homogeneous linear equation iff linear independence of them}
For $n$ $(n-1)$ times differentiable solutions $y_1,y_2,\ldots,y_n$ of $y^{(n)}+\sum_{i=0}^{n-1}p_i(x)y^{(i)}=0$ on an interval $I$ where $p_i(x)$ are continuous for all $i\in[0,n-1]\cap\bbZ$, for any $x_0\in I$, $W(y_1,y_2,\ldots,y_n)(x_0)\ne 0$ iff $y_1,y_2,\ldots,y_n$ are linearly independent.
\sssc{Annihilator}
An annihilator of a function $f$ is a linear differential operator $L$ such that $L(f)=0$.
\sssc{Annihilator multiplication property}
If $P$ is an annihilator of $f$, then for any linear differential operator $Q$, $PQ$ is an annihilator of $f$.
\ssc{Characteristic equation or auxiliary equation for constant coefficient homogeneous linear ODE}
\sssc{Method of solution}
Consider an $n$th-order constant coefficient homogeneous linear ODE
\[\sum_{j=0}^na_jy^{(j)}=0,\quad a_n\neq 0.\]
The formation of its general solutions using the characteristic equation (also called auxiliary equation) is as follows:
\ben
\item The characteristic equation is
\[\sum_{j=0}^na_jr^j=0.\]
\item Suppose there are $k$ distinct real roots $r_1,r_2,\ldots,r_k$, each $r_j$ with multiplicity $m_j$, and $h$ distinct pairs of conjugate imaginary roots $\alpha_1\pm\beta_1i,\alpha_2\pm\beta_2i,\ldots,\alpha_h\pm\beta_hi$ ($\beta_j\neq 0$), each $\alpha_j\pm\beta_ji$ with multiplicity $\omega_j$. (Thus, $\sum_{j=1}^km_j+2\sum_{j=1}^h\omega_j=n$.)
\item The $n$-parameter family of general solutions of the ODE is
\[y=\sum_{j=1}^k\sum_{u=0}^{m_j-1}A_{ju}x^ue^{r_jx}+\sum_{j=1}^h\sum_{u=0}^{\omega_j-1}\qty(B_{ju}\cos(\beta_jx)+C_{ju}\sin(\beta_jx))x^ue^{\alpha_jx},\]
with parameters $A_{ju},B_{ju},C_{ju}$.
\een
\sssc{Exponential order}
For any solution $f(t)$ of any homogeneous linear ODE of one variable $t$, there exists some $M,T\geq 0$ and $\alpha\in\bbR$ such that for all $t>T$, $\abs{f(t)}\leq e^{M\alpha t}$, and the smallest such $\alpha$ is the largest real part among all roots of its characteristic equation.
\ssc{Method of undetermined coefficients (MUC) for constant coefficient non-homogeneous linear ODE}
\sssc{Superposition approach}
Consider an $n$th-order constant coefficient non-homogeneous linear ODE
\[\sum_{j=0}^na_jy^{(j)}=g(x),\quad a_n\neq 0,\]
with $g(x)$ being able to be written in the form
\[g(x)=\sum_{j=1}^pk_jx^{m_j}e^{\alpha_jx}+\sum_{j=p+1}^{p+q}x^{m_j}e^{\alpha_jx}\qty(k_j\cos(\beta_jx)+l_j\sin(\beta_jx)),\]
for some constants $p,q,m_j\in\bbN_0$ and $\alpha_j,\beta_j,k_j,l_j\in\bbR$.

The formation of its general solutions using the method of undetermined coefficients is as follows:
\ben
\item Solve the complementary equation
\[\sum_{j=0}^na_jy^{(j)}=0,\quad a_n\neq 0\]
to obtain its general solutions $y_c$.
\item For each term $k_jx^{m_j}e^{\alpha_jx}$ of $g(x)$ for $j\in\bbN\land j\leq p$, construct a particular solution $y_j$:
\ben
\item If $\alpha_j$ is not a root of the characteristic equation of the complementary equation, let $s_j=0$; otherwise, let $s_j$ be the multiplicity of $\alpha_j$ as a root of the characteristic equation of the complementary equation.
\item The ansatz is
\[y_j=x^{s_j}e^{\alpha_jx}\sum_{k=0}^{m_j}C_{jk}x^k,\]
with constants $C_{jk}\in\bbR$.
\item Solve for $C_{jk}$ such that $y_j$ is a particular solution of the constant coefficient non-homogeneous linear ODE
\[\sum_{j=0}^na_jy^{(j)}=k_jx^{m_j}e^{\alpha_jx}.\]
\een
\item For each term $x^{m_j}e^{\alpha_jx}\qty(k_j\cos(\beta_jx)+l_j\sin(\beta_jx))$ of $g(x)$ for $j\in\bbN\land p<j\leq p+q$, construct a particular solution $y_j$:
\ben
\item If $\alpha_j\pm\beta_ji$ is not a pair of conjugate imaginary roots of the characteristic equation of the complementary equation, let $s_j=0$; otherwise, let $s_j$ be the multiplicity of $\alpha_j\pm\beta_ji$ as a pair of conjugate imaginary roots of the characteristic equation of the complementary equation.
\item The ansatz is
\[y_j=x^{s_j}e^{\alpha_jx}\sum_{k=0}^{m_j}\qty(C_{jk}\cos(\beta_jx)+D_{jk}\sin(\beta_jx))x^k,\]
with constants $C_{jk},D_{jk}\in\bbR$.
\item Solve for $C_{jk},D_{jk}$ such that $y_j$ is a particular solution of the constant coefficient non-homogeneous linear ODE
\[\sum_{j=0}^na_jy^{(j)}=x^{m_j}e^{\alpha_jx}\qty(k_j\cos(\beta_jx)+l_j\sin(\beta_jx)).\]
\een
\item A particular solution of the origin ODE is
\[y_p=\sum_{j=1}^{p+q}y_j,\]
and the general solutions of the origin ODE is
\[y=y_c+y_p.\]
\een
\sssc{Annihilator approach}
Consider an $n$th-order constant coefficient non-homogeneous linear ODE
\[\sum_{j=0}^na_jy^{(j)}=g(x),\quad a_n\neq 0,\]
with $g(x)$ being able to be written in the form
\[g(x)=\sum_{j=1}^pk_jx^{m_j}e^{\alpha_jx}+\sum_{j=p+1}^{p+q}x^{m_j}e^{\alpha_jx}\qty(k_j\cos(\beta_jx)+l_j\sin(\beta_jx)),\]
for some constants $p,q,m_j\in\bbN_0$ and $\alpha_j,\beta_j,k_j,l_j\in\bbR$.

The formation of its general solutions using the method of undetermined coefficients is as follows:
\ben
\item Solve the complementary equation
\[\sum_{j=0}^na_jy^{(j)}=0,\quad a_n\neq 0\]
to obtain its general solutions $y_c$.
\item For each term $k_jx^{m_j}e^{\alpha_jx}$ of $g(x)$, a lowest-order annihilator is $(D-\alpha_j)^{m_j+1}$
\bpr
Prove by Laplace transform.
\[\mathcal{L}\{k_jx^{m_j}e^{\alpha_jx}\}=\frac{k_jm_j!}{(s-\alpha_j)^{m_j+1}}\]
\epr
\item For each term $x^{m_j}e^{\alpha_jx}\qty(k_j\cos(\beta_jx)+l_j\sin(\beta_jx))$ of $g(x)$, a lowest-order annihilator is $(D^2-2\alpha_jD+\alpha_j^{\pht{j}2}+\beta_j^{\pht{j}2})^{m_j+1}$
\bpr
Prove by Laplace transform.
\[\ba
\mathcal{L}\{x^{m_j}e^{\alpha_jx}\qty(k_j\cos(\beta_jx)+l_j\sin(\beta_jx))\}&=\mathcal{L}\{x^{m_j}e^{\alpha_jx}\qty((k_j+l_j)e^{i\beta_jx}-(k_j-l_j)ie^{-i\beta_jx})\}\\
&=(k_j+l_j)\mathcal{L}\{x^{m_j}e^{\alpha_jx+i\beta_jx}\}-(k_j-l_j)i\mathcal{L}\{x^{m_j}e^{\alpha_jx-i\beta_jx}\}\\
&=\frac{(k_j+l_j)m_j!}{(s-\alpha_j-i\beta_j)^{m_j+1}}-\frac{(k_j-l_j)im_j!}{(s-\alpha_j+i\beta_j)^{m_j+1}}
\ea\]
\epr
\item Find a lowest-order annihilator of $g(x)$, $L_1$, and apply it on both sides to obtain a homogeneous linear ODE
\[L_1(\sum_{j=0}^na_jy^{(j)})=0.\]
\item Solve it and obtain its general solutions.
\item Substitute the solutions into the original ODE to obtain an explicit particular solution $y_p$, and the general solutions of the origin ODE is
\[y=y_c+y_p.\]
\een
\ssc{Variation of parameters or variation of constants method of solution of non-homogeneous linear equation}
\sssc{Method of solution}
Given an $n$th-order non-homogeneous linear ODE
\[y^{(n)}+\sum_{i=0}^{n-1}p_i(x)y^{(i)}=q(x),\]
where $p_i(x)$ for all $i\in [0,n-1]\cap\bbZ$ are continuous on An interval $I$.
\ben
\item Let $y_1(x),y_2(x),\ldots,y_n(x)$ be a basis of the vector space of the solutions of the complementary equation,
\[y^{(n)}+\sum_{i=0}^{n-1}p_i(x)y^{(i)}=0.\]
\item The particular solution $y_p$ is given by:
\[y_p=\sum_{i=1}^ny_i(x)\int_{x_0}^x\frac{W_i(t)}{W(t)}\dd{t},\]
where
\bit
\item $x_0$ is an arbitrary point in $I$,
\item $W(x)$ is the Wrońskian of $y_1(x),y_2(x),\ldots,y_n(x)$,
\item $W_i(x)$ is the Wrońskian of $W(x)$ but with the $i$th column replaced by $(0,0,\ldots,0,q(x))$.
\eit
\item The general solutions are given by
\[y=y_c+y_p.\]
\een
\sssc{For second-order equation, Green's function}
For second-order equation
\[y''+p_1(x)y'+p_0(x)y=q(x).\]
Let $y_1(x),y_2(x)$ be two linearly independent solutions of its complementary equation
\[y''+p_1(x)y'+p_0(x)y=0.\]
The particular solution (or family of general solutions with the integral constants being parameters) is given by
\[y_p(x)=-y_1(x)\int_{x_0}^x\frac{y_2(t)q(t)}{y_1(t)y_2'(t)-y_2(t)y_1'(t)}\dd{t}+y_2(x)\int_{x_0}^x\frac{y_1(t)q(t)}{y_1(t)y_2'(t)-y_2(t)y_1'(t)}\dd{t}.\]
This can be written as
\[y_p(x)=\int_{x_0}^xG(x,t)q(t)\dd{t},\]
where
\[G(x,t)=\frac{y_1(t)y_2(x)-y_1(x)y_2(t)}{y_1(t)y_2'(t)-y_2(t)y_1'(t)}\]
and is called the Green's function for the ODE and the linear differential operator $D^2+p_1(x)D+p_0(x)$.
\ssc{Reduction of order}
\sssc{Method}
Given an $n$th-order homogeneous linear ODE
\[\sum_{i=0}^np_i(x)y^{(i)}=0,\]
if a nontrivial particular solution $y_1(x)$ has been known, we can reduce the order of the ODE by 1 by the substitution
\[y=vy_1\]
and then the substitution
\[u=v'\]
to an $(n-1)$th-order linear ODE.
\sssc{Closed-form formula for second-order homogeneous linear ODEs}
Given a second-order homogeneous linear ODE
\[y''+p(x)y'+q(x)y=0,\]
if a nontrivial particular solution $y_1(x)$ that has been known, then the general solutions of the ODE is
\[y=y_1\int\frac{e^{-\int p(x)\dd{x}}}{y_1(x)^2}\dd{x}.\]



\end{document}
